\section{Research scope, status, and conventions}
\label{scope:section}

\subsection{The source state and the questions answered}

The source snapshot is the public ProveIt repository at commit
\begin{center}
\texttt{5a790187c8e186e41e2b990b4941cb7a1a3c7b6b},
\end{center}
inspected on 11 October 2026 (UTC). The two requested directories are
\src{Analysis/Polylogarithms/docs/manuscript} and \src{docs/incoming}
\cite{RepoMain,RepoIncoming}. All eleven incoming archives in that
snapshot were inventoried and extracted. The mathematical review
concentrated on the relevant harmonic, directional-zeta, mixed-spectral,
and Gaussian-certificate sections; this is not a claim to have
independently verified every page of the multi-volume manuscript.

The latest incoming reports already contain substantial results.
The diagonal cubic digamma finite part is reduced to a third diagonal
Tornheim derivative, all quadratic Tornheim directions are known, the
all-nonpositive inner-order sector has a finite rational reduction,
and common-shift harmonic Laurent formulas and Newton interpolations
are available. Repeating those theorems would not address the next
questions. Table~\ref{scope:map} records the precise continuations here.

\begin{table}[htbp]
\centering\small
\begin{tabular}{@{}L{0.28\textwidth}L{0.43\textwidth}L{0.22\textwidth}@{}}
\toprule
Source question & Answer in this article & Status \\
\midrule
Independent harmonic exponents (Exact Identities, Question 8)
& Entire ordered normalization, explicit relative endpoints and
  spectral consequences; Section~\ref{eht:sec}
& Analytic-normalization component answered; full question remains open \\
\addlinespace
All asymmetric cubic Tornheim rays (directional report, Question 2)
& Formula for every admissible direction in two specified scalar
  coordinates; Section~\ref{tr3:section}
& Proved reduction; no independence or ordinary-constant closure claimed \\
\addlinespace
Mixed inner signs at arbitrary depth (directional report, Question 11)
& Finite algorithm retaining a free outer order and repeated colors;
  Section~\ref{mixed:section}
& Proved; exact compiler supplied \\
\addlinespace
Scalar identities from integer Fourier samples (mixed-spectral report)
& Uniqueness on logarithmically dense samples and complete multicentre
  criterion; Section~\ref{sampling:section}
& Proved in the stated finite power--logarithm class \\
\addlinespace
Short Gaussian $S_6$ and revised $S_8$
& Search scope and unchanged conjectural status;
  Section~\ref{audit:section}
& Unresolved \\
\bottomrule
\end{tabular}
\caption{Contribution map. ``Proved'' refers to the stated theorem,
not to a claim of literature-wide priority.}\label{scope:map}
\end{table}

The literature reconciliation is substantive. The one-endpoint
entire interpolant is precisely the function constructed by Ihara,
Nakamura, and Yamamoto \cite{eht:INY}; agreement is established by
their finite reversed-star formula, not by guessing from integer
values. Its identification supplies the independent-order entire
normalization requested in one component of Question 8 of the
exact-identities report \cite{RepoExact}. That question also asks for
a smallest generating kernel with a Gamma-separated singular part;
neither minimality nor that additional factorization is established
here. The additional calculations
are presented as exact consequences and useful research organization.
The mixed-order and Tornheim theorems are advances relative to the
specified source snapshot. The literature search was targeted and
does not establish an absolute priority claim for every formula.

\subsection{What a reduction means}

A reduction is stated relative to a specified output class. A finite
sum of multiple polylogarithms with a free complex first index is an
exact and useful answer even when it does not collapse to ordinary
Gamma and zeta values. Likewise, expressing all cubic directions in
$\Omega$ and $\kappa$ is a theorem about a finite spanning collection;
it is not a theorem that these coordinates are irreducible or
arithmetically independent. In the harmonic sector, the separate
order jets retain the nonsymmetric coordinate $\eta(a)$ from the
ordered-resonance report \cite{RepoOrdered}.

These distinctions prevent three common overstatements: treating a
new name for an integral as its evaluation, treating a finite
normal form as a minimal basis, and treating equality at special
integer parameters as permission to differentiate in a transverse
parameter. Every differentiation below is performed on an identity
valid in a neighborhood in the variable being differentiated.

\subsection{Conventions needed throughout}

The superscript $\zeta^{(j)}(s,a)$ denotes differentiation in the
first, spectral variable. The generalized Stieltjes constants are
fixed by
\begin{equation}\label{scope:Stieltjes}
 \zeta(1+u,a)=\frac1u+
       \sum_{j\ge0}\frac{(-1)^j\gamma_j(a)}{j!}u^j,
 \qquad \gamma_0(a)=-\psi(a).
\end{equation}
Write $\gamma_j=\gamma_j(1)$ and $\gamma=\gamma_0$.
For a finite upper bound, the strict harmonic convention is
\[
 H_N(s_1,\ldots,s_r)
   =\sum_{N\ge n_1>\cdots>n_r\ge1}
                          n_1^{-s_1}\cdots n_r^{-s_r},
 \qquad H_N^{(s)}=H_N(s).
\]
The empty word has value one. The local Hurwitz notation in
Section~\ref{eht:sec} starts its shifted summation at zero, and its
finite-interval formula specifies the conversion explicitly.
Bernoulli numbers use $B_1=-1/2$ and
$te^{xt}/(e^t-1)=\sum_{n\ge0}B_n(x)t^n/n!$.

The notation $\CT_{C=c}$ always means the Laurent constant in the
displayed affine coordinate $C-c$. It does not mean substitution of
individual singular summands. A derivative of a restricted Tornheim
ray means that the ray is formed first and its resulting one-variable
meromorphic germ is then continued and differentiated. Coherent
principal branches are used initially in right half-planes; any
extension to other domains is along specified analytic continuation.
Normalized primitives include their lower endpoint or their value at
the origin, so their additive constants are fixed.

All unrestricted statements are proved analytically or by exact finite
algebra in the article. The accompanying computations test conventions,
coefficients, and independent evaluation routes. Floating-point
agreement is reported as a diagnostic, never as a proof or as an
interval-certified error bound.
